\chapter{Кинематика системы подвижности на кривошипно-шатунных опорах} \label{chapt2}

\section{Вывод уравнений обратной задачи кинематики для платформ с кривошипно-шатунными опорами} \label{sec:InverseKinematics}

\subsection{ОЗК для трехстепенной СП}
\label{subsec:InvKin3DOFMotionPlatform}
Кинематическая схема СП 3DOF показана на рисунке \ref{img:3dof_kinematics_scheme}.
~Введены следующие обозначения:
$R$ -- длина кривошипа, $L$ -- длина шатуна, $a$ -- расстояние по оси X до двигателей 1 и 2, $b$ - расстояние по оси Y до двигателей 1 и 2, $c$ -- расстояние по оси X до двигателя 3, $\alpha$ -- угол крена, $\beta$ -- угол тангажа, $z_p = \left(O О_1\right)$ -- подъем центра подвижной платформы, $h = z_p - L$ -- смещение центра по высоте относительно среднего значения $z_p$, $q_1, q_2, q_3$ -- углы поворота кривошипов.

\begin{figure}[h]
	\center
	\includegraphics[width=0.75\textwidth]{3dof_kinematics_scheme}
	\caption{Кинематическая схема подвижности 3DOF}
	\label{img:3dof_kinematics_scheme}
\end{figure}

В ходе решения ОЗК необходимо определить вектор $\mathbf{q}=(q_1^*, q_2^*, q_3^*)^T$ обобщенных коодинат~по~заданному вектору~положения~платформы $\mathbf{X}=(\alpha,~\beta,~z_p^*)^T$. Для вывода уравнений ОЗК выразим координаты точек С1, С2, С3 через углы крена, тангажа и подъем. Для этого запишем матрицы поворота относительно осей $x$ и $y$

\begin{equation}
	\label{eq:KinematicsRotationMatrices}
	\begin{aligned}	
	& R_x = 
	\left(
		\begin{array}{ccc}
			1 &	0 				 &	0 \\
			0 &	\cos\alpha   &	-\sin\alpha\\
			0 &	\sin\alpha	 &	\cos\alpha
		\end{array}
	\right),
	\\
	& Ry = 
	\left(
		\begin{array}{ccc}
			\cos\beta &	0   &	\sin\beta \\
			0              &	1   & 	0\\
			-\sin\beta &	0	& 	\cos\beta
		\end{array}
	\right).
	\end{aligned}
\end{equation}

Тогда координаты точек крепления шатунов к подвижной платформе:
\begin{equation}
\label{eq:KinematicsUpperPoints}
	\begin{aligned}
		\left(\begin{array}{ccc} x_{C1} \\ y_{C1} \\ z_{C1} \end{array}\right) = R_x R_y
		\left(\begin{array}{ccc} -a \\ -b \\ 0	\end{array}\right) +
		\left(\begin{array}{ccc} 0 \\ 0 \\ z_p^* \end{array}\right) = 
		\left(\begin{array}{lr} 
				a\cos{\beta} \\ 
				-b\cos{\alpha}-a\sin{\alpha}\sin{\beta} \\ 
				z_p^* - b\sin{\alpha} + a\cos{\alpha}\sin{\beta}\end{array}\right)\\
		\left(\begin{array}{ccc} x_{C2} \\ y_{C2} \\ z_{C2} \end{array}\right) = R_x R_y
		\left(\begin{array}{ccc} -a \\ b \\ 0	\end{array}\right) +
		\left(\begin{array}{ccc} 0 \\ 0 \\ z_p^* \end{array}\right) =
		\left(\begin{array}{lr} 
				a\cos{\beta} \\ 
				b\cos{\alpha}-a\sin{\alpha}\sin{\beta} \\ 
				z_p^* + b\sin{\alpha} + a\cos{\alpha}\sin{\beta}\end{array}\right)\\
		\left(\begin{array}{ccc} x_{C3} \\ y_{C3} \\ z_{C3} \end{array}\right) = R_x R_y
		\left(\begin{array}{ccc} c \\ 0 \\ 0	\end{array}\right) +
		\left(\begin{array}{ccc} 0 \\ 0 \\ z_p^* \end{array}\right) = 
		\left(\begin{array}{lr} 
				c\cos{\beta} \\ 
				c\sin{\alpha}\sin{\beta} \\ 
				z_p^* - c\cos{\alpha}\sin{\beta}\end{array}\right)
	\end{aligned}
\end{equation}

Т.к. кривошипы вращаются строго в плоскости $yOz$, справедливы следующие соотношения:
\begin{equation}
\label{eq:KinematicsCrankPoints}
	\begin{aligned}
		& \left(\begin{array}{lr} x_{B1} \\ y_{B1} \\ z_{B1} \end{array}\right) = 
		\left(\begin{array}{lr} -a \\ -b + R\sin{q_1^*} \\ -R\cos{q_1^*} \end{array}\right);\\
		& \left(\begin{array}{lr} x_{B2} \\ y_{B2} \\ z_{B2} \end{array}\right) = 
		\left(\begin{array}{lr} -a \\ b + R\sin{q_2^*} \\ -R\cos{q_2^*} \end{array}\right);\\
		& \left(\begin{array}{lr} x_{B3} \\ y_{B3} \\ z_{B3} \end{array}\right) = 
		\left(\begin{array}{lr} c \\ R\sin{q_3^*} \\ -R\cos{q_3^*} \end{array}\right).
	\end{aligned}
\end{equation}

Далее, т.к. расстояние между точками $B_i$ и $C_i$, $i=1\ldots3$ неизменно и равно $L$, можно записать следующие уравнения:
\begin{equation}
\label{eq:KinematicsRodLengths}
	\begin{aligned}
		 (x_{B1} - x_{C1})^2 + (y_{B1} - y_{C1})^2 + (z_{B1} - z_{C1})^2 = L^2;\\
		 (x_{B2} - x_{C2})^2 + (y_{B2} - y_{C2})^2 + (z_{B2} - z_{C2})^2 = L^2;\\
		 (x_{B3} - x_{C3})^2 + (y_{B3} - y_{C3})^2 + (z_{B3} - z_{C3})^2 = L^2.
	\end{aligned}
\end{equation}

После подстановки \ref{eq:KinematicsCrankPoints} и \ref{eq:KinematicsUpperPoints} в \ref{eq:KinematicsRodLengths} и дополнительных преобразований получаем следующие уравнения:

\begin{equation}
\label{eq:KinematicsDriveAnglesFromPosEq}
\begin{aligned}
\bar{A_i} + \bar{B_i}\sin{q_i} + \bar{C_i}\cos{q_i} = 0,~~i = 1,2,3
\end{aligned}
\end{equation}

где коэффициенты $\bar{A_i}, \bar{B_i}, \bar{C_i}$ зависят от $\alpha, \beta, z_p$ в соответствии с формулами \ref{eq:KinematicsUpperPoints},\ref{eq:KinematicsCrankPoints}.

\begin{equation}
\label{eq:KinematicsSquareEqKoefficients}
\begin{aligned}
	& \bar{A_i} = (x_{Ci} - x_{Ai})^2 + (y_{Ci} - y_{Ai})^2 + (z_{Ci} - z_{Ai})^2 + R^2 + L^2,\\
	& \bar{B_i} = -2R(y_{Ci} - y_{Ai}),\\
	& \bar{C_i} = 2R(z_{Ci} - z_{Ai}).
\end{aligned}
\end{equation}

Перейдя от $\sin{q_i}$ и $\cos{q_i}$ к половинному углу через $\tg\left({q_i/2}\right)$, получим следующие квадратные уравнения:
\begin{equation}
\label{eq:KinematicsSquareEquations}
\begin{aligned}
	\left(\bar{A_i}-\bar{C_i}\right)\tg^2{\frac{q_i}{2}} + 2B_i\tg{\frac{q_i}{2}} + \left(\bar{A_i}+\bar{C_i}\right)=0
\end{aligned}
\end{equation}
Из решения уравнения \ref{eq:KinematicsSquareEquations} получаем два корня:

\begin{equation}
\label{eq:KinematicsPhiRoots}
\begin{aligned}
	q_i = -2\arctan\frac{\bar{B_i}\pm\sqrt{\bar{B_i^2} + \bar{C_i^2} - \bar{A_i^2}}}{\bar{A_i} - \bar{C_i}}
\end{aligned}
\end{equation}

Оба корня с точки зрения кинематики платформы идентичны. Выбором знака <<+>> или <<-->> в уравнении \ref{eq:KinematicsPhiRoots} можно задать область значений обобщенной координаты $(q_i < 0)$ или $(q_i > 0)$, т.е. определить соответствующую рабочую полуплоскость вращения кривошипа. Как будет показано далее, выбор рабочей плоскости сказывается на динамических характеристиках СП.

\subsection{ОЗК для двухстепенной СП}

Система подвижности с двумя степенями свободы, изображенная на рисунке \ref{img:2dof_kinematics_scheme} представляет собой частный случай рассмотренной выше кинематической схемы 3DOF.

\begin{figure}
	\center
	\includegraphics[width=0.75\textwidth]{2dof_kinematics_scheme}
	\caption{Кинематическая схема подвижности 2DOF}
	\label{img:2dof_kinematics_scheme}
\end{figure}

Координаты точек $B_1,~B_2,~C_1,~C_2$ определяются по формулам \ref{eq:KinematicsUpperPoints}.

Получить уравнения ОЗК для данного типа СП можно путем задания константы $z_p$ и исключения третьего уравнения из \ref{eq:KinematicsRodLengths}.
Тогда результат решения ОЗК для случая 2DOF имеет вид:

\begin{equation}
\label{eq:Kinematics2DOFResult}
\begin{aligned}
	& q_i = -2\arctan\frac{\bar{B_i}\pm\sqrt{\bar{B_i^2} + \bar{C_i^2} - \bar{A_i^2}}}{\bar{A_i} - \bar{C_i}},~i = 1,2\\
	& \bar{A_i} = (x_{Ci} - x_{Ai})^2 + (y_{Ci} - y_{Ai})^2 + (z_{Ci} - z_{Ai})^2 + R^2 + L^2,\\
	& \bar{B_i} = -2R(y_{Ci} - y_{Ai}),\\
	& \bar{C_i} = 2R(z_{Ci} - z_{Ai}).
\end{aligned}
\end{equation}
\subsection{ОЗК для одностепенной СП}

Рассмотрим кинематику платформы с одной степенью свободы, схема которой изображена на рис \ref{img:1dof_kinematics_scheme}

\begin{figure}
	\center
	\includegraphics[width=0.9\textwidth]{1dof_kinematics_scheme}
	\caption{Кинематическая схема подвижности 1DOF}
	\label{img:1dof_kinematics_scheme}
\end{figure}

Данная кинематическая схема позволяет осуществлять движение закрепленного на подвижном основании объекта лишь по оси тангажа. Для решения задачи управления платформой необходимо определить угол поворота кривошипа $q$ по заданному углу тангажа $\beta$. ОЗК для данного типа СП решается аналогично рассмотренному подходу в \ref{subsec:InvKin3DOFMotionPlatform}. Уравнение для длины шатуна в случае одностепенной СП имеет вид:

\begin{equation}
\label{eq:Kinematics1DOFRodLengths}
\begin{aligned}
	& \left(x_B-x_A\right)^2 + \left(y_B-y_A\right)^2 = L^2,\\
	& x_B = \sqrt{L1^2 + L2^2}\cos{\psi} - c,~y_B = \sqrt{L1^2 + L2^2}\sin{\psi} + a,\\
	& x_A = R\cos{\psi},~y_A = R\sin{\psi},\\
	& \psi = \pi+\beta-\arctan{\frac{L_2}{L_1}}.
\end{aligned}
\end{equation}

Упростив уравнение \ref{eq:Kinematics1DOFRodLengths}, и проведя аналогичные \ref{eq:KinematicsSquareEquations} преобразования, получим следующее решение ОЗК для 1DOF СП:

\begin{equation}
\label{eq:Kinematics1DOFSolution}
\begin{aligned}
	& q = -2\arctan\frac{\bar{B}\pm\sqrt{\bar{B^2} + \bar{C^2} - \bar{A^2}}}{\bar{A} - \bar{C}}, где \\
	& \bar{A} = L_1^2 + L_2^2 + R^2 + a^2 + c^2 - L^2 + 2\sqrt{L_1^2 + L_2^2}\left(ccos{\psi}-a\sin{\psi}\right),\\
	& \bar{B} = -2R\left(\sqrt{L_1^2 + L_2^2}\sin{\psi}\right) - a),\\
	& \bar{C} = -2R\left(\sqrt{L_1^2 + L_2^2}\cos{\psi} + c\right).
\end{aligned}
\end{equation}

\section{Прямая задача кинематики для платформ с кривошипно-шатунными опорами}
Решение прямой задачи кинематики позволяет определить положение(ориентацию) подвижной платформы в пространстве по заданным углам поворота кривошипов, что может быть использовано при:
\begin{itemize}
\item определении расхождения между желаемым и действительным положением платформы без использования дополнительных датчиков;
\item построении кинематической ограничительной характеристики;
\item выводе уравнений динамики платформы
\end{itemize}
Получим решение ПЗК на примере трехстепенной динамической платформы.

\subsection{Решение уравнений прямой задачи кинематики численными методами}
Путем подстановки \ref{eq:KinematicsCrankPoints} и \ref{eq:KinematicsUpperPoints} в \ref{eq:KinematicsRodLengths} получим следующие соотношения:
\begin{equation}
\label{eq:DirectKinematicsProblemFormulation}
\begin{aligned}
& F_1\left(X\right)=\left(b\sin{\alpha}sin{\beta}-a\cos{\beta}+a\right)^2+\left(b\cos{\alpha}+b-R\sin{q_1}\right)^2+\\
&+\left(z_p+a\sin{\beta}+b\cos{\beta}\sin{\alpha}+R\cos{q_1}\right)^2-L^2=0;\\
& F_2\left(X\right)=\left(-b\sin{\alpha}sin{\beta}-acos{\beta}+a\right)^2+\left(b\cos{\alpha}+b+R\sin{q_2}\right)^2+\\
&+\left(z_p+a\sin{\beta}-b\cos{\beta}\sin{\alpha}+R\cos{q_2}\right)^2-L^2;\\
& F_3\left(X\right)=\left(c\cos{\beta}-c\right)^2+\left(R\sin{q_3}\right)^2+\left(z_p-c\sin{\beta}+R\cos{q_3}\right)^2-L^2=0.
\end{aligned}
\end{equation}

Таким образом, требуется определить $X$, решив систему нелинейных уравнений $F\left(X\right)=\mathbf{0}$. Применим рекуррентную формулу метода Ньютона для систем линейных алгебраических уравнений:
\begin{equation}
\label{eq:DirectKinematicsNewton}
\begin{aligned}
	& J \left( X^k\right) \Delta X^k = -F\left(X^k\right);\\
	& X^{k+1} = X^k + \Delta X^k,~k=1,2,\dots,N;\\
	& \left|X^{k+1}-X^k\right|<\varepsilon,
\end{aligned},
\end{equation}

где $J\left( X^k\right)$ -- матрица Якоби системы $F\left( X^k\right)=\mathbf{0}$:
\begin{equation}
\label{eq:DirectKinematicsJacobi}
\begin{aligned}
	J\left( X^k\right)=\left(
	\begin{array}{ccc}
		\frac{F_1\left(x_1^k\right)-F_1\left(x_1^{k-1}\right)}{x_1^k-x_1^{k-1}} & 
		\frac{F_1\left(x_2^k\right)-F_1\left(x_2^{k-1}\right)}{x_2^k-x_2^{k-1}} & 
		\frac{F_1\left(x_3^k\right)-F_1\left(x_3^{k-1}\right)}{x_3^k-x_3^{k-1}} \\
		\frac{F_2\left(x_1^k\right)-F_2\left(x_1^{k-1}\right)}{x_1^k-x_1^{k-1}} & 
		\frac{F_2\left(x_2^k\right)-F_2\left(x_2^{k-1}\right)}{x_2^k-x_2^{k-1}} & 
		\frac{F_2\left(x_3^k\right)-F_2\left(x_3^{k-1}\right)}{x_3^k-x_3^{k-1}} \\
		\frac{F_3\left(x_1^k\right)-F_3\left(x_1^{k-1}\right)}{x_1^k-x_1^{k-1}} & 
		\frac{F_3\left(x_1^k\right)-F_3\left(x_2^{k-1}\right)}{x_2^k-x_2^{k-1}} & 
		\frac{F_3\left(x_1^k\right)-F_3\left(x_3^{k-1}\right)}{x_3^k-x_3^{k-1}}
	\end{array}		
	\right)
\end{aligned}
\end{equation}
\subsection{Аппроксимация уравнений ПЗК}

Для дальнейшего решения и анализа уравнений динамики СП численного решения ПЗК оказывается недостаточно, вследствие чего требуется также получить аналитические зависимости $\mathbf{X}=f\left(\mathbf{q}\right)$. 

Т.к. уравнения системы \ref{eq:DirectKinematicsProblemFormulation} напрямую не обращаются, необходимо упростить выражения, в которые входят тригонометрические и степенные функции от $\alpha,~\beta,~h$. Для этого разложим $\alpha, \beta, h$ уравнения  в ряд по степеням $\kappa = R/L$: 
\begin{equation}
\label{eq:DirectKinematicsSeriesExpansion}
\begin{aligned}
	\alpha = C_{1,\alpha}\kappa + C_{2,\alpha}\kappa^2, \\
	\beta = C_{1,\beta}\kappa + C_{2,\beta}\kappa^2, \\
	h = C_{1,h}\kappa + C_{2,h}\kappa^2.
\end{aligned}
\end{equation}

Для разложения выбрана величина $R/L$, т.к. вследствие особенностей проектирования КШМ, данное соотношение всегда меньше единицы и обычно находится в пределах 0.1-0.2. Следовательно, высокими порядками $\kappa$ можно пренебречь без существенных потерь точности вычислений $\alpha,~\beta,~h$. Ограничимся многочленами второй степени в \ref{eq:DirectKinematicsSeriesExpansion}.
	
Для получения коэффициентов $C_{1,\alpha},C_{2,\alpha},C_{1,\beta},C_{2,\beta}$ раскроем уравнения длин шатунов \ref{eq:KinematicsRodLengths}:
\begin{equation}
\label{eq:DirectKinematicsRodLengthsFull}
	\begin{aligned}
	0 = & h^2 + 2hL + 2a^2\left(1-\cos\beta\right) + 2b^2\left(1-\cos\alpha\right) + R^2 + 2Rb\sin{q_1}\left(1 - \right.\\
	&\left.-\cos\alpha\right) + 2\left(h+L\right)R\cos{q_1} + 2\left(h+L\right)\left(a\sin\beta + b\sin\alpha\cos\beta\right)+\\
	& + 2b\sin\alpha\left(a\sin\beta + R\cos{q_1}\cos\beta\right) + 2aR\cos{q_1}\sin\beta \\
	0 = & h^2 + 2hL + 2a^2\left(1-\cos\beta\right) + 2b^2\left(1-\cos\alpha\right) + R^2 - 2Rb\sin{q_2}\left(1 -\right.\\ 
	&\left. - \cos\alpha\right) + 2\left(h+L\right)R\cos{q_2} + 2\left(h+L\right)\left(a\sin\beta - b\sin\alpha\cos\beta\right)-\\
	& - 2b\sin\alpha\left(a\sin\beta + R\cos{q_2}\cos\beta\right) + 2aR\cos{q_2}\sin\beta\\
	0 = & h^2 + 2hL + 2c^2\left(1-\cos\beta\right) + R^2 + 2\left(h+L\right)R\cos{q_3}-\\
	& - 2\left(h+L\right)c\sin\beta - 2cR\cos{q_3}\sin\beta.
	\end{aligned}
\end{equation}
	
Для тригонометрических функций применим разложение Тейлора того же порядка:
\begin{equation}
\label{eq:DirectKinematicsSinCosSeries}
\begin{aligned}
	& \sin\alpha \approx \alpha-\frac{\alpha^3}{6} \approx C_{1,\alpha}\kappa + C_{2,\alpha}\kappa^2,\\
	& \cos\alpha \approx 1 - \frac{\alpha^2}{2} \approx 1 - \frac{C^2_{1,\alpha}}{2}\kappa^2,\\
	& \sin\beta \approx \beta-\frac{\beta^3}{6} \approx C_{1,\beta}\kappa + C_{2,\beta}\kappa^2,\\
	& \cos\beta \approx 1 - \frac{\beta^2}{2} \approx 1 - \frac{C^2_{1,\beta}}{2}\kappa^2.
\end{aligned}
\end{equation}	

В результате решения полиномиальной относительно $\kappa$ системы получим следующие коэффициенты разложения функций $\alpha,\beta,h$.

Коэффициенты первого порядка при $\kappa$:
\begin{equation}
\label{eq:DirectKinematicsCoefsApprox1}
\begin{aligned}
	& C_{1,\alpha} = \frac{L}{2b}\left(\cos{q_2} - \cos{q_1}\right),\\
	& C_{1,\beta} = \frac{L}{a+c}\left(\cos{q_3} - \frac{1}{2}\left(\cos{q_1} + \cos{q_2}\right)\right),\\
	& C_{1,h} = -\frac{L}{2\left(a+c\right)}\left(2a\cos{q_3} + c\left(\cos{q_1} + \cos{q_2}\right)\right).
\end{aligned}
\end{equation}

Коэффициенты второго порядка при $\kappa^2$:
\begin{equation}
\label{eq:DirectKinematicsCoefsApprox2}
\begin{aligned}
	C_{2,\alpha} & = - \frac{1}{2}\left(\cos{q_1} + \cos{q_2}\right)C_{1,\alpha}=\frac{L}{4b}\left(\cos^2{q_2} - \cos^2{q_1}\right)\\
	C_{2,\beta} & = -\frac{L}{2\left(a+c\right)}\left[\cos^2{q_3} - \frac{1}{2}\left(\cos^2{q_1} + \cos^2{q_2}\right)\right],\\
	C_{2,h} & = cC_{2,\beta} + \frac{L}{2}\left(\cos^2{q_3} - 1\right) = \\
	& = \frac{L}{2\left(a+c\right)}\left[a\cos^2{q_3} + \frac{c}{2}\left(\cos^2{q_1} + \cos^2{q_2}\right)\right] - \frac{L}{2}.
\end{aligned}
\end{equation}

Окончательно, подставив \ref{eq:DirectKinematicsCoefsApprox1} и \ref{eq:DirectKinematicsCoefsApprox2} в \ref{eq:DirectKinematicsSeriesExpansion} получим следующие формулы для $\alpha,~\beta,~h$:
\begin{equation}
\label{eq:DireKinematicsAnalytic}
\begin{aligned}
	& h \approx \frac{-R}{1+c/a}\left[\cos{q_3} + \frac{c}{2a}\left(\cos{q_1} + \cos{q_2}\right) + \frac{R}{2L}\left(\sin^2{q_3} + \right.\right.\\
	& \left.\left. + \frac{c}{2a}\left(\sin^2{q_1} + \sin^2{q_2}\right)\right)\right], \\
	& \alpha \approx \frac{R}{2b}\left[\cos{q_2}-\cos{q_1} + \frac{R}{2L}\left(\cos^2{q_1}-\cos^2{q_2}\right)\right],\\
	& \beta \approx \frac{R}{a+c}\left[\cos{q_3} - \frac{1}{2}\left(\cos{q_1} + \cos{q_2}\right)\right].
\end{aligned}
\end{equation}

Для определения точности формул \ref{eq:DireKinematicsAnalytic} зададим равномерную сетку значений $\left(\alpha,~\beta,~h\right)$ и вычислим соответствующие комбинации $\left(q_1,~q_2,~q_3\right)$ с помощью ОЗК. Далее по полученным приближенным формулам вычисляются оценки $(\hat{\alpha},~\hat{\beta},~\hat{h})$. Максимальные расхождения между значениями $\left(\alpha,~\beta,~h\right)$ и $(\hat{\alpha},~\hat{\beta},~\hat{h})$ для нескольких СП с различным $\kappa$ приведены в таблице \ref{tbl:DirectKinematicsMaxErrors}.

\begin{table} [htbp]%
    \centering
	\caption{Ошибка приближения прямой задачи кинематики для различных СП}%
	\label{tbl:DirectKinematicsMaxErrors}
    \renewcommand{\arraystretch}{1.5}
    \begin{SingleSpace}
	\begin{tabular}{@{}@{\extracolsep{20pt}}lllll@{}} 
        \toprule     %%% верхняя линейка
    	п/п	& $\kappa$ &$E_{max, \alpha}, ^o$	& $E_{max, \beta}, ^o$& $E_{max, h}$, мм	\\
        \midrule
    	СП1 	& 0.28 	 & 0.26		& 0.16	&	0.39\\
    	СП2		& 0.20 	 & 0.66		& 0.68 &		0.85\\
    	СП3 	& 0.27 	 & 0.69		& 0.53 &		1.46\\
    	СП4     & 0.20 	 & 0.21		& 0.29	&	0.29\\
        \bottomrule
	\end{tabular}%
   	\end{SingleSpace}
\end{table}

\begin{table} [htbp]%
    \centering
	\caption{Параметры СП }%
	\label{tbl:DirectKinematicsMaxErrorsMPParams}
    \renewcommand{\arraystretch}{1.5}
    \begin{SingleSpace}
	\begin{tabular}{@{}@{\extracolsep{20pt}}llllll@{}} 
        \toprule     %%% верхняя линейка
    	п/п	& $a$, м & $b$, м & $c$, м & $R$, м & $L$, м \\
        \midrule
    	СП1 & 0.122 & 0.195 & 0.250 & 0.045 & 0.160 \\
    	СП2 & 0.170 & 0.290 & 0.340 & 0.100 & 0.500	\\
    	СП3 & 0.170 & 0.400 & 0.560 & 0.130 & 0.560	\\
    	СП4 & 0.170 & 0.300 & 0.340 & 0.070 & 0.430	\\
        \bottomrule
	\end{tabular}%
   	\end{SingleSpace}
\end{table}

Под данным таблицы \ref{tbl:DirectKinematicsMaxErrors} видно, что максимальная погрешность не превышает $0.25^o$ для углов крена и тангажа и 0.1 мм для подъема, что позволяет использовать формулу \ref{eq:DireKinematicsAnalytic} для дальнейших расчетов.
Рассмотрим распределение ошибки прямого преобразования координат в пределах рабочей зоны платформы. Т.к. полные 3D-фигуры плохо читаются, на рисунках \ref{img:DirectKinematicErrors} приведены срезы рабочей зоны координатными плоскостями $\alpha O \beta$, $\alpha O h$, $\beta O h$.

%Ошибка приближения по высоте
\begin{figure}[h!]
    \center{
        \hfill
        \subcaptionbox[List-of-Figures entry]{\label{img:DirKinematicsErrorHwithAlphaFixed}} {\includegraphics[width=0.5\linewidth]{dir_kin_error_h_with_alpha0}}%
        \hfill       
        \subcaptionbox{\label{img:DirKinematicsErrorHwithBetaFixed}} {\includegraphics[width=0.5\linewidth]{dir_kin_error_h_with_beta0}}
        \hfill
    }
    \center{
        \subcaptionbox[List-of-Figures entry]{\label{img:DirKinematicsErrorAlphaWithHFixed}} {\includegraphics[width=0.5\linewidth]{dir_kin_error_alpha_with_h0}}%
        \hfill       
        \subcaptionbox{\label{img:DirKinematicsErrorAlphaWithBetaFixed}} {\includegraphics[width=0.5\linewidth]{dir_kin_error_alpha_with_beta0}}
        \hfill
    }
	\center{    
        \subcaptionbox[List-of-Figures entry]{\label{img:DirKinematicsErrorBetaWithHFixed}} {\includegraphics[width=0.5\linewidth]{dir_kin_error_beta_with_h0}}%
        \hfill       
        \subcaptionbox{\label{img:DirKinematicsErrorBetaWithAlphaFixed}} {\includegraphics[width=0.5\linewidth]{dir_kin_error_beta_with_alpha0}}
        \hfill
    }
    \caption{Ошибка приближенного решения ПЗК} % Этот текст попадает в названия рисунков в списке рисунков
    \label{img:DirectKinematicErrors}
\end{figure}

%%Ошибка приближения по углу крена
%\begin{figure}[ht]
%    \center{
%        \hfill
%        
%    }
%    \caption{Ошибка приближенного решения ПЗК по $\alpha$} % Этот текст попадает в названия рисунков в списке рисунков
%    \label{img:DirectKinematicErrorsAlpha}
%\end{figure}
%
%%Ошибка приближения по углу тангажа
%\begin{figure}[ht]
%    \center{
%        \hfill
%        
%    }
%    \caption{Ошибка приближенного решения ПЗК по $\beta$} % Этот текст попадает в названия рисунков в списке рисунков
%    \label{img:DirectKinematicErrorsBeta}
%\end{figure}

\section{Расчет векторов шатунов}
Для решения задач динамики необходимо определять положение точек шатунов в каждый момент времени.

Пусть заданы углы $q_1,~q_2,~q_3$. Векторы вдоль шатунов от кривошипов к платформе обозначим $S_1,~S_2,~S_3$. 

С точностью до третьего порядка относительно $\kappa$ их компоненты определяются следующим образом.
	
Первый шатун:

\begin{equation}
\begin{aligned}
\frac{1}{L}S_1 &= \begin{pmatrix}0 \\ 0 \\ 1\end{pmatrix} + \kappa\begin{pmatrix}0 \\ -\sin{q_1}\\ 0\end{pmatrix} + \kappa^2\begin{pmatrix}\frac{a}{2L}C^2_{1,\beta} + \frac{b}{L}C_{1,\alpha}^2\\-\frac{b}{2L}C_{1,\alpha}^2\\ -\frac{1}{2}\sin^2{q_1} \\ \end{pmatrix} + \\
& + \kappa^3\frac{1}{L}\begin{pmatrix}aC_{1,\beta}C_{2,\beta} + bC_{1,\alpha}C_{2,\beta} + bC_{2,\alpha}C_{1,\beta} \\ -bC_{1,\alpha}C_{2,\alpha} \\ -\frac{b}{2}\sin{q_1}C_{1,\alpha}^2\end{pmatrix}
\end{aligned}
\end{equation}

Для второго шатуна:

\begin{equation}
\begin{aligned}
\frac{1}{L}S_2 &= \begin{pmatrix}0 \\ 0 \\ 1\end{pmatrix} + \kappa\begin{pmatrix}0 \\ -\sin{q_2}\\ 0\end{pmatrix} + \kappa^2\begin{pmatrix}\frac{a}{2L}C^2_{1,\beta} + \frac{b}{L}C_{1,\alpha}^2\\\frac{b}{L}C_{1,\alpha}C_{1,\beta} -\frac{b}{2L}C_{1,\alpha}^2\\ -\frac{1}{2}\sin^2{q_1} \\ \end{pmatrix} + \\
& + \kappa^3\frac{1}{L}\begin{pmatrix}aC_{1,\beta}C_{2,\beta} + bC_{1,\alpha}C_{2,\beta} + bC_{2,\alpha}C_{1,\beta} \\ -bC_{1,\alpha}C_{2,\alpha} \\ -\frac{b}{2}\sin{q_1}C_{1,\alpha}^2\end{pmatrix}
\end{aligned}
\end{equation}

Третий шатун:

\begin{equation}
\begin{aligned}
	\frac{1}{L}S_3 = \begin{pmatrix}0 \\ 0 \\ 1\end{pmatrix} + \kappa\begin{pmatrix}0 \\ -\sin{q_3}\\ 0\end{pmatrix} + \kappa^2\begin{pmatrix}-\frac{c}{2L}C^2_{1,\beta}\\0\\-\frac{1}{2}\sin^2{q_3}\end{pmatrix} + \kappa^3\frac{1}{L}\begin{pmatrix}-\frac{c}{L}C_{1,\beta}C_{2,\beta}\\0\\0\end{pmatrix}
\end{aligned}
\end{equation}

\subsection{Расчет векторов шатунов в системе координат платформы}

Для дальнейших расчетов также потребуется вычисление координат точек шатунов в подвижной системе координат.

Пересчет координат осуществляется по формуле:
\begin{equation}
\begin{aligned}
	r^0 = R_xR_y\left(r - e_3\left(L + h\right)\right)
\end{aligned}
\end{equation}

Для сокращения последующих выкладок введем обозначения для модифицированных синусов:

\begin{equation}
\begin{aligned}
	sq_i = \sin{q_i} - C_{1,\alpha},\quad i =1,\ldots,3
\end{aligned}
\end{equation}

для показателей крена и тангажа:

\begin{equation}
\begin{aligned}
	\mu & = \frac{1}{4}C_{1,\alpha}\left(\sin{q_2} - \sin{q_1}\right) = \frac{L}{8b}\left(\cos{q_2}-\cos{q_1}\left(\sin{q_2}-\sin{q_1}\right)\right),\\
	\lambda & = \frac{1}{4}\left(2\sin{q_3}-\sin{q_1}-\sin{q_2}\right)C_{1,\alpha}
\end{aligned}
\end{equation},

скорректированного второго коэффициента $C_{2,\beta}$:

\begin{equation}
\begin{aligned}
	\gamma_1 & = \frac{a}{2L}C^2_{1,\beta} - C_{2,\beta} + \frac{b}{L}C_{1,\alpha}C_{1,\beta}, \\ 
	\gamma_2 & = \frac{a}{2L}C^2_{1,\beta} - C_{2,\beta} - \frac{b}{L}C_{1,\alpha}C_{1,\beta}, \\
	\gamma_3 & = -\frac{c}{2L}C^2_{1,\beta} - C_{2,\beta}.
\end{aligned}
\end{equation}

для скорректированного второго коэффициента -- $C_{2,\alpha}$:

\begin{equation}
\begin{aligned}
	\delta_1 = C_{2,\alpha} - \frac{b}{2L}C^2_{1,\alpha},\quad C_{2,\alpha} + \frac{b}{2L}C^2_{1,\alpha},\quad \delta_3 = C_{2,\alpha}.
\end{aligned}
\end{equation}

Тогда для первого шатуна:

\begin{equation}
\begin{aligned}
	\frac{1}{L}S_1^0 = \begin{pmatrix} 0 \\ 0 \\ 1\end{pmatrix} + \kappa\begin{pmatrix}-C_{1,\beta} \\ -sq_1 \\ 0 \end{pmatrix} + \kappa^2\begin{pmatrix}\gamma_1 \\ \delta_1 \\ -\frac{1}{2}\left(C^2_{1,\beta} + sq^2_1\right)\end{pmatrix} + \\ \kappa^3\begin{pmatrix}
	-\mu \frac{b}{a+c}C_{1,\alpha} + \frac{b}{L}C_{1,\alpha}C_{2,\beta} - \frac{1}{L}C_{1,\beta}C_{2,h}\\
	\frac{a}{L}C_{1,\alpha}C_{2,\beta} + \frac{1}{L}C_{1,\alpha}C_{2,h} - \mu C_{1,\alpha} \\ 
	C_{1,\beta}\gamma_1 + sq_1\delta_1
	\end{pmatrix}
\end{aligned}
\end{equation}

Для второго шатуна:

\begin{equation}
\begin{aligned}
	\frac{1}{L}S_1^0 = \begin{pmatrix} 0 \\ 0 \\ 1\end{pmatrix} + \kappa\begin{pmatrix}-C_{1,\beta} \\ -sq_1 \\ 0 \end{pmatrix} + \kappa^2\begin{pmatrix}\gamma_1 \\ \delta_1 \\ -\frac{1}{2}\left(C^2_{1,\beta} + sq^2_1\right)\end{pmatrix} + \\ \kappa^3\begin{pmatrix}
	-\mu \frac{b}{a+c}C_{1,\alpha} + \frac{b}{L}C_{1,\alpha}C_{2,\beta} - \frac{1}{L}C_{1,\beta}C_{2,h}\\
	\frac{a}{L}C_{1,\alpha}C_{2,\beta} + \frac{1}{L}C_{1,\alpha}C_{2,h} - \mu C_{1,\alpha} \\ 
	C_{1,\beta}\gamma_1 + sq_1\delta_1
	\end{pmatrix}
\end{aligned}
\end{equation}

Для третьего шатуна:

\begin{equation}
\begin{aligned}
	\frac{1}{L}S_1^0 = \begin{pmatrix} 0 \\ 0 \\ 1\end{pmatrix} + \kappa\begin{pmatrix}-C_{1,\beta} \\ -sq_1 \\ 0 \end{pmatrix} + \kappa^2\begin{pmatrix}\gamma_1 \\ \delta_1 \\ -\frac{1}{2}\left(C^2_{1,\beta} + sq^2_1\right)\end{pmatrix} + \\ \kappa^3\begin{pmatrix}
	-\mu \frac{b}{a+c}C_{1,\alpha} + \frac{b}{L}C_{1,\alpha}C_{2,\beta} - \frac{1}{L}C_{1,\beta}C_{2,h}\\
	\frac{a}{L}C_{1,\alpha}C_{2,\beta} + \frac{1}{L}C_{1,\alpha}C_{2,h} - \mu C_{1,\alpha} \\ 
	C_{1,\beta}\gamma_1 + sq_1\delta_1
	\end{pmatrix}
\end{aligned}
\end{equation}

\section{Ограничительная кинематическая характеристика СП}
Расчет ограничительной кинематической характеристики СП -- поверхности, построенной по границе рабочей зоны платформы, производится с целью планирования траекторий движения и определения предельных характеристик платформы по углам крена, тангажа и высоте. 

Пусть $X_{up}$ -- точка крепления одного из шатунов к платформе, а $X_e$ -- центр соответствующего вала. Введем обозначения:

\begin{itemize}
\item $\Delta X = X_{up} - X_e = \left(\Delta x,~\Delta y,~\Delta z\right)^T$ -- вектор, соединяющий центра вала с точкой крепления шатуна с платформой;
\item $L_yz = \sqrt{L^2 - \Delta x^2}$ -- длина проекции шатуна на плоскость $OYZ$;
\item $\left|\Delta X_{yz}\right| = \sqrt{\Delta y^2 + \Delta z^2}$ -- длина проекции вектора $\Delta X$ на плоскость $OYZ$.

Тогда для сушествования угла поворота $q_i$ кривошипа, реализующего точку $X_up$ необходимо и достаточно выполнения неравенств:

\begin{equation}
\label{eq:AngleRealizableEquation}
\begin{aligned}
	\left|L_{yz}-R\right| \leq \| \Delta X_{yz} \| \leq L_{yz} + R
\end{aligned}
\end{equation}
\end{itemize}

Фиксируем $\alpha, \beta$. Выделим в $\Delta z = \Delta z^0 + h$ слагаемое $\Delta z^0$, не зависящее от $h$. Пусть выполнено условие $L_{yz} - R \geqslant \left|\Delta y\right|$.

Тогда неравенства \ref{eq:AngleRealizableEquation} равносильны неравенствам:

\begin{equation}
\begin{aligned}
-\Delta z^0 + \sqrt{\left(L_{yz}-R\right)^2 - \Delta y^2} \leq h \leq -\Delta z^0 + \sqrt{\left(L_{yz}+R\right)^2 - \Delta y^2}
\end{aligned}
\end{equation}

Главная часть в ограничениях -- это первое слагаемое. При небольших углах $\alpha, \beta$ линейная аппроксимация дает:

\begin{equation}
\begin{aligned}
-\Delta z_1^0 & \approx -\alpha\beta - b\alpha, \\
-\Delta z_2^0 & \approx -\alpha\beta + b\alpha, \\
-\Delta z_3^0 & \approx c\beta. \\
\end{aligned}
\end{equation}

Поэтому приближенные формулы принимают вид:
\begin{equation}
\begin{aligned}
-a\beta + b\alpha - R \leq & h \leq -a\beta + b\alpha + R, \\
c\beta - R \leq h \leq c\beta + R.
\end{aligned}
\end{equation}

Пересечение непусто, если пара углов $\left(\alpha,\beta\right)$ удовлетворяет ограничениям:
\begin{equation}
\begin{aligned}
	b\left|\alpha\right| & \leq R \\
	b\left|\alpha\right| + \left(a+c\right)\left|\beta\right| & \leq 2R
\end{aligned}
\end{equation}

которые определяют шестиугольник с вершинами:
\begin{equation}
\begin{aligned}
	\begin{pmatrix}0 & \pm\frac{2R}{a+c}\end{pmatrix}, ~\begin{pmatrix}\pm\frac{R}{b} & \pm\frac{R}{a+c}\end{pmatrix} 
\end{aligned}
\end{equation}

Пусть фиксирован уровень $h$. Тогда множество углов $\left(\alpha,~\beta\right)$, совместимых с этим уровнем, приближенно определяется неравенствами:
\begin{equation}
\begin{aligned}
	b\left|\alpha\right| + \left|a\beta + h\right| &\leq R, \\
	\left|c\beta - h\right| \leq R.
\end{aligned}
\end{equation}

Это множество может быть четырехугольником, пятиугольником или шестиугольником. Оно получается пересечением ромба с вершинами в точках:
\begin{equation}
\begin{aligned}
	\begin{pmatrix}0 & -\frac{h \pm R}{a}\end{pmatrix}, ~\begin{pmatrix}\pm\frac{R}{b} & -\frac{h}{a}\end{pmatrix} 
\end{aligned}
\end{equation}

с полосой
\begin{equation}
\begin{aligned}
	\frac{h-R}{c} \leq \beta \frac{h+R}{c}
\end{aligned}
\end{equation}

Это пересечение непусто тогда и только тогда, когда:
\begin{equation}
\begin{aligned}
	-R \leq z \leq R
\end{aligned}
\end{equation}

Окончательно, для трехстепенной платформы максимальные углы крена, тангажа и максимальный подъем определяются по следующим формулам: 

\begin{equation}
\label{eq:KinematicConstraints}
\begin{aligned}
	& \left|\alpha\right|\leq\alpha_{max}=\frac{R}{b},\\
	& \left|\beta\right|\leq\beta_{max}=\frac{2R}{a+c},\\
	& \left|h\right| \leq h_{max}=R.
\end{aligned}
\end{equation}

\begin{figure}[h!]
    \center{
        \hfill
        \subcaptionbox[List-of-Figures entry]{\label{img:RestrictKinematicAlphaHNS}} {\includegraphics[width=0.5\linewidth]{RestrictKinematicAlphaHNS}}%
        \hfill       
        \subcaptionbox{\label{img:RestrictKinematicAlphaH}} {\includegraphics[width=0.5\linewidth]{RestrictKinematicAlphaH}}
        \hfill
    }
    \center{
        \subcaptionbox[List-of-Figures entry]{\label{img:RestrictKinematicBetaHNS}} {\includegraphics[width=0.5\linewidth]{RestrictKinematicBetaHNS}}%
        \hfill       
        \subcaptionbox{\label{img:RestrictKinematicBetaH}} {\includegraphics[width=0.5\linewidth]{RestrictKinematicBetaH}}
        \hfill
    }
	\center{    
        \subcaptionbox[List-of-Figures entry]{\label{img:RestrictKinematicAlphaBetaNS}} {\includegraphics[width=0.5\linewidth]{RestrictKinematicAlphaBetaNS}}%
        \hfill       
        \subcaptionbox{\label{img:RestrictKinematicAlphaBeta}} {\includegraphics[width=0.5\linewidth]{RestrictKinematicAlphaBeta}}
        \hfill
    }
    \caption{Проекции зоны кинематических ограничений для симметричной (а,в,д) и несимметричной (б,г,е) СП} % Этот текст попадает в названия рисунков в списке рисунков
    \label{img:RestrictKinematicZone}
\end{figure}


